\section{Locating the point}\label{sec:point}

The preceding section states what membership would require: a point
$(p,q)\in U_0$ and a $t\in\mathbb C^{*}$ at which the family
\eqref{eq:fam-a}--\eqref{eq:fam-c} reproduces System~$\mathsf C$'s projective
normal form \eqref{eq:rC}. This section supplies such a point, checks the
hypotheses under which the moduli space is built, and then shows that the two
solutions found are \emph{all} of them. What it does not do is remove the word
\emph{projective}; that is \S\ref{sec:projective}.

\subsection{The dictionary, and the entry that failed}\label{ssec:dictionary}

Matching \eqref{eq:rC} against the family requires knowing how the family's
parameters $(\theta_0,\theta_1,\theta_\infty,t)$ appear in the projective
normal potential $R$. The dictionary was not assumed; it was established by
probing three exact draws of the family and reading off the correspondence
(\textbf{VERIFIED}, \texttt{S16-08}). Three entries were confirmed on all three
draws: the coefficient of the double pole at $z=0$ is $(\theta_0^2-1)/4$, the
coefficient of the double pole at $z=1$ is $(\theta_1^2-1)/4$, and
$R(\infty)=t^2/4$.

A fourth entry \emph{failed}, and the failure is reported here because the
located value of $\theta_\infty$ is unintelligible without it. The naive
reading of the $1/z$ coefficient at infinity, $t\,\theta_\infty/2$, was
rejected on all three draws; the draws give
\[
  \frac{t\,(\theta_\infty+1)}{2},
\]
the $+1$ being the shift introduced by passing to a scalar reduction. That
shift is the entire reason the located exponent is
$\theta_\infty=-2\mp\sqrt3/9$ rather than $-1\mp\sqrt3/9$. A reader who
carries the naive entry will find the value off by exactly one and will look
for the error in the wrong place.

The three draws were chosen to carry $\theta_0\neq0$, which makes them
simultaneously the vacuity control for the hypothesis check of
\S\ref{ssec:hyp-check}: the observable does report a nonzero $\theta_0$ when
one is present, so its report of $\theta_0=0$ for System~$\mathsf C$ is
informative and not an artefact.

\subsection{The three hypotheses at System \texorpdfstring{$\mathsf
C$}{C}}\label{ssec:hyp-check}

\paragraph{Normalised singular points (\texttt{S16-07}, \textbf{PROVEN}).}
Under the M\"obius map $x=s_++\delta z$ of \eqref{eq:spm} the potential
\eqref{eq:rC} has poles only at $z=0$ and $z=1$, both double with coefficient
$-\tfrac14$, and no apparent singularity (\texttt{S16-06}, recorded at
\S\ref{ssec:sysBC}). At infinity the transported potential expands as
\[
  r_z=-\frac{11}{27}+\frac{i\sqrt{11}/27}{z}+\cdots,
\]
so the normalisation $t^2/4=R(\infty)$ gives
\[
  t=\mp\frac{2i\sqrt{33}}{9}.
\]
This is \emph{exactly} the value the companion record prints for
System~$\mathsf C$. The companion's $t$ is therefore van der Put--Saito's time
invariant, computed from the configuration; there is no mismatch between the
two to explain. The agreement is with an exact algebraic number, not a
numerical coincidence, which is why the label is \textbf{PROVEN}.

\paragraph{Irreducibility (\texttt{S16-10}, \textbf{PROVEN}).}
System~$\mathsf C$ is irreducible over $\mathbb C(x)$, proved exactly by the
Riccati test. A solution of $w'+w^2=r$ can have only simple poles; the residues
at $s_\pm$ satisfy $\kappa^2-\kappa=-\tfrac14$, so $\kappa=\tfrac12$; $w$ is bounded at
infinity with $c_0^2=\tfrac13$; and matching the $1/x$ coefficient of $r$,
which is $-\tfrac19$, forces
\[
  n=\frac{-1/9}{2c_0}-1=-1\mp\frac{\sqrt3}{18},
\]
which is irrational and hence not a non-negative integer. The required
contradiction follows, and System~$\mathsf C$ does not sit on the classical
Riccati locus.

\begin{remark}[A void check, retained; \texttt{S16-11},
\textbf{PROVEN}]\label{rem:void}
The first irreducibility script reached the verdict \emph{irreducible} with no
test performed. Symbolic cancellation of $(x-c)^2r$ evaluated at an algebraic
$c$ returned \texttt{nan}, so the list of admissible residues came out empty,
the product over that empty list was trivially the required value, and a
companion test of the form ``every multiplicity equals~$2$'' was vacuously true
on an empty dictionary. The defect was caught only because a
reducible-by-construction control was \emph{also} declared irreducible. The
void script is retained in the record. The general rule it forced --- a step
that cannot fail cannot be trusted to report failure, so every control must be
an assertion --- is applied throughout this paper's supporting computations.
\end{remark}

\paragraph{Local data (\texttt{S16-09}, \textbf{STRUCTURAL}).}
System~$\mathsf C$ has $\theta_0=\theta_1=0$, hence $\det A_i=0$ and
$\operatorname{tr}A_i=0$, so $A_0$ and $A_1$ are nilpotent and nonzero: the
local data are non-semisimple, and \S\ref{ssec:nonapparent} already records the
logarithmic behaviour this produces. This was \emph{contrary to
pre-registration}, which recorded the expectation that the hypothesis would be
violated or left undecided. It is not violated, for two reasons. The eigenvalue
difference is $0$, which is not in $\mathbb Z\setminus\{0\}$, so the data are
not resonant in van der Put--Saito's sense; and, as \S\ref{ssec:resonance}
records, no semisimplicity condition is imposed in the construction in any
case. Van der Put--Saito admit precisely the logarithmic local data that
Jimbo's~(A1) excludes.

The label here is \textbf{STRUCTURAL} rather than \textbf{PROVEN}, and the
reason is worth stating plainly: the conclusion rests on the \emph{absence} of
a hypothesis in the source (\texttt{S16-01}, \textbf{INSPECTED}). An absence
read from a page is weaker evidence than a computation, and no amount of
re-reading converts it.

\subsection{The located point}\label{ssec:located}

With the dictionary of \S\ref{ssec:dictionary} in hand, the match is an
overdetermined algebraic system: five coefficient equations from matching two
quartics in $z$, together with the no-apparent-singularity condition $b_1=0$,
in the four unknowns $p,q,t,\theta_\infty$. Here $b_1$ is the coefficient of
$z$ in the entry $b$ of \eqref{eq:fam-b}, which is affine in $z$; so $b_1=0$
says exactly that $b$ is constant, and \S\ref{ssec:tracefree} is where that
matters.
It is consistent, and it has a
solution on each branch of $t$ (\textbf{VERIFIED}, \texttt{S16-12}).

On the branch $t=-2i\sqrt{33}/9$,
\[
  \theta_\infty=-2-\frac{\sqrt3}{9},\qquad
  q=\frac12+\frac{i\bigl(-457937\sqrt{11}+281700\sqrt{33}\bigr)}{122408}
   \approx 0.5+0.8123\,i,
\]
\[
  p=\frac{5i\bigl(-28151758571\sqrt{33}+48874838016\sqrt{11}\bigr)}
         {4086468672}\approx 0.4649\,i .
\]
On the branch $t=+2i\sqrt{33}/9$,
\[
  \theta_\infty=-2+\frac{\sqrt3}{9},\qquad
  q\approx 0.5-25.6278\,i,\qquad p\approx 396.2089\,i .
\]

The values were produced by a hand-reduced closed form, and are therefore
checked by a second route (\textbf{VERIFIED}, \texttt{S16-13}): re-substituting
each solution into the \emph{general} formula for $R$ --- not the reduced form
used to derive it --- makes $R-r_z$ simplify identically to zero, with worst
relative residual over five sample points at $60$ digits of
$5.79877\times10^{-56}$ on the first branch and $7.00553\times10^{-61}$ on the
second. In both cases $b_1=0$ and $q\notin\{0,1\}$, so each point lies in the
chart $U_0$ of \eqref{eq:chart}.

Two things this second route does not do. It is a re-derivation by a different
route, not an independent confirmation by a different party: the same agent
wrote both. And it certifies the values it was given; it cannot certify that
there are no others. That is the next subsection.

\subsection{The location is exhaustive}\label{ssec:exhaustive}

The two solutions of \S\ref{ssec:located}, one on each branch of $t$, would be
an empty inventory if
the solver that produced the two solutions had silently missed a third. It did
not, and this is settled exactly rather than by trusting the solver
(\textbf{PROVEN}, \texttt{S25-06}).

\paragraph{Transcription controls (\texttt{S25-01}, \textbf{PROVEN}).}
The six generators of the matching system were re-derived independently of the
script that first solved it. A positive control: both recorded solutions
annihilate all six generators exactly --- as symbolic zero, not as numerical
smallness --- and both satisfy $q(q-1)\neq0$. A negative control: perturbing
$\theta_\infty$ by $1/1000$ at the first solution makes generators $1$, $2$,
$3$ and $5$ nonzero, so the generator set is not identically zero. Both arms
were pre-registered and run before the main result was consulted.

\paragraph{The base field, and why it is not larger.}
The radicals visible above are $\sqrt{33}$, $\sqrt{11}$ and $\sqrt3$, which
invites the field $\mathbb Q(\sqrt3,\sqrt{11})$. The correct base field is
\[
  K=\mathbb Q\bigl(i\sqrt{11}\bigr),
\]
because that is the field in which the matching equations' \emph{own}
coefficients live, by way of $s_+=-\tfrac16+\tfrac{i\sqrt{11}}{6}$ and
$\delta=-\tfrac{i\sqrt{11}}{3}$ in \eqref{eq:spm}. Adjoining $\sqrt3$ in
advance would not be a harmless convenience: $\sqrt3$ is \emph{generated} by
the answer, and pre-adjoining it would split the solution set into two
rational points and make a degree-$2$ ideal look like two components of
degree~$1$ --- destroying exactly the statement the computation exists to make.

\paragraph{The ideal (\texttt{S25-02}, \textbf{PROVEN}).}
The ideal generated by the six matching generators, saturated at $q(q-1)\neq0$
by the Rabinowitsch variable $u$ with $u\,q(q-1)-1=0$, is zero-dimensional of
degree exactly~$2$ over $K$. It is computed as a count of standard monomials
under a \texttt{grevlex} Gr\"obner basis over an exact algebraic number field,
so the result is exact. The lexicographic eliminant in $\theta_\infty$ is
\begin{equation}\label{eq:eliminant}
  27\,\theta_\infty^2+108\,\theta_\infty+107=0,
  \qquad
  \theta_\infty=-2\pm\frac{\sqrt3}{9},
\end{equation}
and \eqref{eq:eliminant} was registered as a prediction \emph{before} the
computation was run.

\begin{remark}[Galois conjugacy, and what it is not]\label{rem:galois}
The two branches of \S\ref{ssec:located} are therefore exactly the Galois
conjugates of one another over $K$ under $\sqrt3\mapsto-\sqrt3$, exhibited by
the eliminant \eqref{eq:eliminant} rather than observed numerically. This is
the same symmetry that relates the two branches at the level of the character
variety (\texttt{S22-20}, \textbf{VERIFIED}), where the two branches give the
same point up to an interchange of $x_1$ and $x_2$.

It does \emph{not} close the $z$-dependent gauge question, which is an open
question of this paper (\texttt{S17-07}, \textbf{CONJECTURED}) and is stated as
one in \S\ref{sec:open}. Two statements in this paper carry a $\sqrt3$; they
are unrelated, and a reader should not infer a connection from the coincidence
of a radical.
\end{remark}

\paragraph{One of the six conditions was already implied (\texttt{S25-03},
\textbf{PROVEN}).}
The no-apparent-singularity condition $b_1=0$ lies in the radical of the ideal
generated by the five coefficient-matching equations together with the
Rabinowitsch generator --- indeed in the ideal itself, since $b_1$ reduces to
$0$ modulo the \texttt{grevlex} Gr\"obner basis, and independently since
adjoining $1-v\,b_1$ yields the unit ideal. So matching the two quartics
\emph{already forces} $b_1=0$: the system of \S\ref{ssec:located} is
overdetermined by two, not by one. The membership test carries its own control
--- $\theta_\infty$, which vanishes nowhere on the variety by
\eqref{eq:eliminant}, is correctly reported as \emph{not} in the radical --- so
the test is not accepting everything put to it.

\paragraph{Closure, by two routes (\texttt{S25-06}, \textbf{PROVEN}).}
The solution set, saturated at $q(q-1)\neq0$, is exactly the two points of
\S\ref{ssec:located}. Counting route: the degree is $2$, which bounds the
number of points over the algebraic closure counted with multiplicity by $2$,
while the positive control exhibits two \emph{distinct} points, so each has
multiplicity~$1$ and there is no room for a third. Constructive route: the
lexicographic Gr\"obner basis in the order $u>p>q>t>\theta_\infty$ is
triangular, with $u$, $p$, $q$ and $t$ each appearing to degree exactly~$1$ in
a basis element whose only other variable is $\theta_\infty$; so each root of
\eqref{eq:eliminant} determines at most one point, the eliminant has exactly
two distinct roots, and back-substitution returns exactly one point per root,
matching the recorded solutions exactly.

The scope of this result is narrow and should be read narrowly: it settles
whether the solver's two solutions were all of them. It says nothing about the
derivation that produced the six generators, which is unchanged here.

\subsection{Membership, and the word it leaves standing}
\label{ssec:membership-verdict}

Collecting: System~$\mathsf C$ satisfies every hypothesis under which van der
Put--Saito build the moduli space for the $(0,0,1,-)$ row --- local data
(\texttt{S16-09}), irreducibility (\texttt{S16-10}), and the normalisation at
infinity (\texttt{S16-07}) --- and it is a point of the chart $U_0$
(\texttt{S16-12}, \texttt{S16-13}), uniquely so up to the Galois conjugate
(\texttt{S25-06}). Since the foliation given by \eqref{eq:omega} coincides with
the isomonodromy foliation (\texttt{S16-03}, \textbf{INSPECTED}), and since
$q(q-1)\neq0$ and $t\neq0$ at the located point so that the vector field
\eqref{eq:flow} is regular there, exactly one Painlev\'e~V leaf passes through
System~$\mathsf C$. The deformation exists; it is not an ansatz.

This conclusion is \textbf{STRUCTURAL} (\texttt{S16-14}), and the caution is
part of the claim rather than a hedge attached to it. What has been matched is
the \emph{projective} normal form. What is established is therefore that
$\mathsf C$'s projectivisation is that of the family member at $(p,q,t)$ ---
equality up to a rank-$1$ twist, in the same sense in which $\mathsf B$ and
$\mathsf C$ are projectively equivalent at \eqref{eq:gauge} without being the
same operator.

Two further limits, stated here so that \S\ref{sec:projective} is read as
removing one of them and not both. The pre-registered monodromy-and-Stokes
falsifier at $30$ digits was \emph{not} run; nothing above rests on it and
nothing above stands in for it. And the gauge question of \S\ref{sec:open}
remains open: the completeness result of \S\ref{ssec:exhaustive} concerns the
solution set of the matching system, not the uniqueness of a representative
within a projective class.
